\chapter*{Словарь терминов}
\addcontentsline{toc}{chapter}{Словарь терминов}

\textbf{auth-jwt} --- проектируемый микросервис авторизации платформы Mindbox,
отвечающий за централизованный выпуск JWT-токенов и ротацию ключей подписи.

\textbf{auth-jwt-operator} --- Kubernetes-оператор системы авторизации (язык Go),
доставляющий токены и открытые ключи (JWKS) от auth-jwt до микросервисов
через Kubernetes Secret и избавляющий их от прямых обращений к auth-jwt.

\textbf{Claim (утверждение)} --- пара <<ключ-значение>> в полезной нагрузке JWT,
описывающая субъект, его права доступа или метаданные токена.
Зарезервированные claim (\texttt{iss}, \texttt{sub}, \texttt{aud}, \texttt{exp}, \texttt{iat})
определены стандартом RFC~7519.

\textbf{Custom Resource (CR)} --- пользовательский тип ресурса,
расширяющий API Kubernetes и описывающий прикладную сущность в декларативном виде;
обрабатывается соответствующим оператором.

\textbf{DREAD} --- модель количественной оценки рисков по пяти критериям
(Damage, Reproducibility, Exploitability, Affected~users, Discoverability),
дополняющая качественную классификацию STRIDE числовым рейтингом.

\textbf{Grace-период} --- временной интервал после ротации ключа подписи,
в течение которого предыдущий ключ остаётся доступным для верификации
ранее выпущенных токенов.

\textbf{Identity Provider (IdP)} --- система управления идентификацией,
отвечающая за аутентификацию пользователей и сервисов
и выпуск токенов безопасности.

\textbf{JSON Web Key (JWK)} --- формат представления криптографических ключей
в виде JSON-объектов, определённый стандартом RFC~7517.

\textbf{JSON Web Token (JWT)} --- открытый стандарт (RFC~7519),
определяющий компактный формат передачи утверждений (claims)
между двумя сторонами в виде подписанного JSON-объекта.

\textbf{JWK Set (JWKS)} --- набор открытых ключей в формате JWK,
публикуемый эмитентом через специальный endpoint.
Потребители получают из него актуальные ключи для проверки подписи,
что обеспечивает автоматическое обнаружение нового ключа при ротации.

\textbf{Kubernetes (K8s)} --- платформа оркестрации контейнеров с открытым исходным кодом,
обеспечивающая автоматическое развёртывание, масштабирование
и управление контейнеризированными приложениями.

\textbf{Kubernetes-оператор} --- приложение, расширяющее Kubernetes:
отслеживает Custom Resource и приводит фактическое состояние кластера
к описанному в спецификации через цикл реконсиляции.
В данной работе доставляет токены и открытые ключи микросервисам
через Kubernetes Secret.

\textbf{Mindbox.Authentication} --- клиентская библиотека платформы Mindbox,
предоставляющая микросервисам абстракции для получения и проверки JWT-токенов.

\textbf{mTLS (mutual TLS)} --- взаимная аутентификация на транспортном уровне,
при которой обе стороны соединения предъявляют сертификаты,
подтверждая свою идентичность.

\textbf{OAuth~2.0} --- открытый стандарт авторизации (RFC~6749),
описывающий делегирование доступа к ресурсам без передачи учётных данных;
служит основой для протоколов OpenID Connect и Token Exchange.

\textbf{OpenID Connect (OIDC)} --- протокол аутентификации,
построенный как слой идентификации поверх OAuth~2.0.

\textbf{Policy Decision Point (PDP)} --- компонент системы авторизации,
принимающий решение о допуске или отклонении запроса
на основе политик доступа.

\textbf{Policy Enforcement Point (PEP)} --- компонент системы авторизации,
перехватывающий запрос и применяющий решение, полученное от PDP.

\textbf{RPS (requests per second)} --- число запросов в секунду;
метрика интенсивности нагрузки на сервис.

\textbf{Service Mesh} --- выделенный инфраструктурный слой
для управления межсервисным взаимодействием в микросервисной архитектуре,
реализуемый через sidecar-прокси.

\textbf{Sidecar} --- вспомогательный контейнер,
развёрнутый рядом с основным микросервисом в одном поде Kubernetes
и перехватывающий его сетевой трафик.

\textbf{SIEM (Security Information and Event Management)} --- класс систем
для централизованного сбора, корреляции и анализа событий безопасности.
В данной работе значимые события авторизации доставляются в SIEM для аудита.

\textbf{SLA (Service Level Agreement)} --- соглашение об уровне обслуживания,
фиксирующее обязательства поставщика перед потребителем
по доступности и производительности сервиса.

\textbf{SLO (Service Level Objective)} --- целевой показатель качества сервиса,
определяющий конкретное числовое значение метрики
(например, задержка на 95-м процентиле не более 30~мс),
достижение которого контролируется командой разработки.

\textbf{STRIDE} --- методология моделирования угроз, разработанная в Microsoft.
Классифицирует угрозы по шести категориям: подмена идентичности (Spoofing),
подделка данных (Tampering), отказ от авторства (Repudiation),
раскрытие информации (Information disclosure), отказ в обслуживании (Denial of service),
повышение привилегий (Elevation of privilege).

\textbf{Terraform Provider (authjwt)} --- плагин Terraform (язык Go),
реализующий CRUD-операции с ресурсами auth-jwt
(клиенты авторизации, доверенные эмитенты) через его API;
позволяет описывать конфигурацию авторизации в подходе IaC.

\textbf{Terraform-модуль} --- переиспользуемый пакет конфигурации (язык HCL),
описывающий всех клиентов авторизации платформы;
единый источник истины о конфигурации, применяемый ко всем контурам
для обеспечения их однородности.

\textbf{Token Exchange} --- протокол (RFC~8693), позволяющий обменять
один токен безопасности на другой у службы авторизации.

\textbf{Зона доступности (availability zone)} --- изолированный сегмент
инфраструктуры провайдера с независимыми электропитанием, охлаждением и сетью;
размещение реплик сервиса в разных зонах обеспечивает устойчивость к отказу целой зоны.

\textbf{Инфраструктура как код (IaC)} --- подход к управлению инфраструктурой,
при котором конфигурация описывается в виде декларативных файлов,
хранящихся в системе контроля версий.

\textbf{Контур} --- в терминологии Mindbox: окружение, представляющее собой
совокупность Kubernetes-кластеров, на которых развёрнут полный набор
микросервисов платформы.

\textbf{Мультитенантная архитектура} --- архитектура программной системы,
при которой один экземпляр приложения обслуживает нескольких клиентов (тенантов),
разделяя данные на уровне приложения.

\textbf{Перцентиль} --- значение, ниже которого лежит заданная доля наблюдений.
Например, 95-й перцентиль (p95) времени ответа~--- порог, который не превышают
95\,\% запросов; используется при формулировке SLO.

\textbf{Реконсиляция (reconciliation)} --- периодический цикл сверки,
в котором Kubernetes-оператор сравнивает желаемое состояние (из Custom Resource)
с фактическим и устраняет расхождения.

\textbf{Ротация ключей} --- процедура плановой или экстренной замены
криптографических ключей подписи с сохранением работоспособности системы.

\textbf{Тенант} --- клиент мультитенантной платформы,
данные и конфигурация которого логически изолированы от других клиентов.

\textbf{Эмитент (issuer)} --- система, выпускающая подписанные токены
и публикующая свои открытые ключи для их верификации.
В контексте данной работы: Kubernetes, GitLab CI, Octopus Deploy, Google Workspace
и другие экземпляры auth-jwt.
